\noindent
\begin{minipage}[t]{0.485\textwidth}
\textbf{10.} \textbf{(a)} Có điều gì sai trong phương trình sau đây?
\[
\frac{x^2 + x - 6}{x - 2} = x + 3
\]
\textbf{(b)} Dựa vào phần (a), hãy giải thích tại sao phương trình
\[
\lim_{x \to 2} \frac{x^2 + x - 6}{x - 2} = \lim_{x \to 2} (x + 3)
\]
lại đúng.

\vspace{6pt}
{\color{stewartcyan}\textbf{11--34}}\enspace Tính giới hạn, nếu tồn tại.

\vspace{4pt}
\noindent
\begin{tabular}{@{}p{0.48\linewidth}@{\hspace{0.04\linewidth}}p{0.48\linewidth}@{}}
\textbf{11.} $\lim\limits_{x \to -2} (3x - 7)$ & \textbf{12.} $\lim\limits_{x \to 6} \left(8 - \dfrac{1}{2}x\right)$ \\[11pt]
\textbf{13.} $\lim\limits_{t \to 4} \dfrac{t^2 - 2t - 8}{t - 4}$ & \textbf{14.} $\lim\limits_{x \to -3} \dfrac{x^2 + 3x}{x^2 - x - 12}$ \\[13pt]
\textbf{15.} $\lim\limits_{x \to -2} \dfrac{x^2 + 5x + 4}{x - 2}$ & \textbf{16.} $\lim\limits_{x \to 4} \dfrac{x^2 + 3x}{x^2 - x - 12}$ \\[13pt]
\textbf{17.} $\lim\limits_{x \to -2} \dfrac{x^2 - x - 6}{3x^2 + 5x - 2}$ & \textbf{18.} $\lim\limits_{x \to -5} \dfrac{2x^2 + 9x - 5}{x^2 - 25}$ \\[13pt]
\textbf{19.} $\lim\limits_{t \to 3} \dfrac{t^3 - 27}{t^2 - 9}$ & \textbf{20.} $\lim\limits_{u \to -1} \dfrac{u + 1}{u^3 + 1}$ \\[13pt]
\textbf{21.} $\lim\limits_{h \to 0} \dfrac{(h - 3)^2 - 9}{h}$ & \textbf{22.} $\lim\limits_{x \to 9} \dfrac{9 - x}{3 - \sqrt{x}}$ \\[13pt]
\textbf{23.} $\lim\limits_{h \to 0} \dfrac{\sqrt{9 + h} - 3}{h}$ & \textbf{24.} $\lim\limits_{x \to 2} \dfrac{2 - x}{\sqrt{x + 2} - 2}$ \\[13pt]
\textbf{25.} $\lim\limits_{x \to 3} \dfrac{\dfrac{1}{x} - \dfrac{1}{3}}{x - 3}$ & \textbf{26.} $\lim\limits_{h \to 0} \dfrac{(-2 + h)^{-1} + 2^{-1}}{h}$ \\[15pt]
\textbf{27.} $\lim\limits_{t \to 0} \dfrac{\sqrt{1 + t} - \sqrt{1 - t}}{t}$ & \textbf{28.} $\lim\limits_{t \to 0} \left( \dfrac{1}{t} - \dfrac{1}{t^2 + t} \right)$ \\[13pt]
\textbf{29.} $\lim\limits_{x \to 16} \dfrac{4 - \sqrt{x}}{16x - x^2}$ & \textbf{30.} $\lim\limits_{x \to 2} \dfrac{x^2 - 4x + 4}{x^4 - 3x^2 - 4}$ \\[13pt]
\textbf{31.} $\lim\limits_{t \to 0} \left[ \dfrac{1}{t\sqrt{1 + t}} - \dfrac{1}{t} \right]$ & \textbf{32.} $\lim\limits_{x \to -4} \dfrac{\sqrt{x^2 + 9} - 5}{x + 4}$ \\[15pt]
\textbf{33.} $\lim\limits_{h \to 0} \dfrac{(x + h)^3 - x^3}{h}$ & \textbf{34.} $\lim\limits_{h \to 0} \dfrac{\dfrac{1}{(x + h)^2} - \dfrac{1}{x^2}}{h}$
\end{tabular}

\vspace{8pt}
\noindent{\color{stewartcyan}\rule{\linewidth}{0.6pt}}
\vspace{6pt}

\graphcalcicon\ \textbf{35.} \textbf{(a)} Ước lượng giá trị của
\[
\lim_{x \to 0} \frac{x}{\sqrt{1 + 3x} - 1}
\]
bằng cách vẽ đồ thị hàm số $f(x) = x/(\sqrt{1 + 3x} - 1)$.\\
\textbf{(b)} Lập bảng các giá trị của $f(x)$ với các giá trị $x$ tiến gần đến $0$ và dự đoán giá trị của giới hạn.\\
\textbf{(c)} Sử dụng các Quy tắc Giới hạn để chứng minh dự đoán của bạn là đúng.
\end{minipage}%
\hfill
\begin{minipage}[t]{0.485\textwidth}
\graphcalcicon\ \textbf{36.} \textbf{(a)} Sử dụng đồ thị của
\[
f(x) = \frac{\sqrt{3 + x} - \sqrt{3}}{x}
\]
để ước lượng giá trị của $\lim\limits_{x \to 0} f(x)$ chính xác đến hai chữ số thập phân.\\
\textbf{(b)} Lập bảng các giá trị của $f(x)$ để ước lượng giới hạn chính xác đến bốn chữ số thập phân.\\
\textbf{(c)} Sử dụng các Quy tắc Giới hạn để tìm giá trị chính xác của giới hạn.

\vspace{6pt}
\graphcalcicon\ \textbf{37.} Sử dụng Định lý Kẹp để chỉ ra rằng
\[
\lim_{x \to 0} x^2 \cos 20\pi x = 0
\]
Minh họa bằng cách vẽ đồ thị các hàm số $f(x) = -x^2$, $g(x) = x^2 \cos 20\pi x$, và $h(x) = x^2$ trên cùng một màn hình.

\vspace{6pt}
\graphcalcicon\ \textbf{38.} Sử dụng Định lý Kẹp để chỉ ra rằng
\[
\lim_{x \to 0} \sqrt{x^3 + x^2} \sin \frac{\pi}{x} = 0
\]
Minh họa bằng cách vẽ đồ thị các hàm số $f$, $g$, và $h$ (theo ký hiệu của Định lý Kẹp) trên cùng một màn hình.

\vspace{6pt}
\textbf{39.} Nếu $4x - 9 \le f(x) \le x^2 - 4x + 7$ với mọi $x \ge 0$, hãy tìm $\lim\limits_{x \to 4} f(x)$.

\vspace{6pt}
\textbf{40.} Nếu $2x \le g(x) \le x^4 - x^2 + 2$ với mọi $x$, hãy tính $\lim\limits_{x \to 1} g(x)$.

\vspace{6pt}
\textbf{41.} Chứng minh rằng $\lim\limits_{x \to 0} x^4 \cos \dfrac{2}{x} = 0$.

\vspace{6pt}
\textbf{42.} Chứng minh rằng $\lim\limits_{x \to 0^+} \sqrt{x}\, e^{\sin(\pi/x)} = 0$.

\vspace{6pt}
{\color{stewartcyan}\textbf{43--48}}\enspace Tìm giới hạn, nếu tồn tại. Nếu giới hạn không tồn tại, hãy giải thích lý do.

\vspace{4pt}
\noindent
\begin{tabular}{@{}p{0.48\linewidth}@{\hspace{0.04\linewidth}}p{0.48\linewidth}@{}}
\textbf{43.} $\lim\limits_{x \to -4} (|x + 4| - 2x)$ & \textbf{44.} $\lim\limits_{x \to -4} \dfrac{|x + 4|}{2x + 8}$ \\[12pt]
\textbf{45.} $\lim\limits_{x \to 0.5^-} \dfrac{2x - 1}{|2x^3 - x^2|}$ & \textbf{46.} $\lim\limits_{x \to -2} \dfrac{2 - |x|}{2 + x}$ \\[14pt]
\textbf{47.} $\lim\limits_{x \to 0^-} \left( \dfrac{1}{x} - \dfrac{1}{|x|} \right)$ & \textbf{48.} $\lim\limits_{x \to 0^+} \left( \dfrac{1}{x} - \dfrac{1}{|x|} \right)$
\end{tabular}

\vspace{8pt}
\noindent{\color{stewartcyan}\rule{\linewidth}{0.6pt}}
\vspace{6pt}

\textbf{49. Hàm Signum}\enspace Hàm signum (hay hàm dấu), ký hiệu bởi $\operatorname{sgn}$, được định nghĩa bởi
\[
\operatorname{sgn} x = \begin{cases}
-1 & \text{nếu } x < 0 \\
\phantom{-}0 & \text{nếu } x = 0 \\
\phantom{-}1 & \text{nếu } x > 0
\end{cases}
\]
\textbf{(a)} Vẽ phác thảo đồ thị của hàm số này.\\
\textbf{(b)} Tìm mỗi giới hạn sau đây hoặc giải thích tại sao không tồn tại.\\[4pt]
\noindent
\begin{tabular}{@{}p{0.48\linewidth}@{\hspace{0.04\linewidth}}p{0.48\linewidth}@{}}
\textbf{(i)} $\lim\limits_{x \to 0^+} \operatorname{sgn} x$ & \textbf{(ii)} $\lim\limits_{x \to 0^-} \operatorname{sgn} x$ \\[8pt]
\textbf{(iii)} $\lim\limits_{x \to 0} \operatorname{sgn} x$ & \textbf{(iv)} $\lim\limits_{x \to 0} |\!\operatorname{sgn} x|$
\end{tabular}
\end{minipage}
